\documentclass[12pt]{amsart}
\usepackage{lmodern}
\usepackage{amsmath,amssymb,amsthm}
\usepackage{hyperref}
\usepackage{enumitem}

\newtheorem{theorem}{Theorem}[section]
\newtheorem{lemma}[theorem]{Lemma}
\newtheorem{proposition}[theorem]{Proposition}
\newtheorem{corollary}[theorem]{Corollary}
\newtheorem{definition}[theorem]{Definition}
\newtheorem{remark}[theorem]{Remark}

\newcommand{\re}{\operatorname{Re}}
\newcommand{\im}{\operatorname{Im}}
\newcommand{\RR}{\mathbb{R}}
\newcommand{\ZZ}{\mathbb{Z}}
\newcommand{\CC}{\mathbb{C}}
\newcommand{\Csub}{\mathcal{C}_{\mathrm{logpoly}}^{\mathrm{sub,ell}}}

\title[Spectral Exclusions for Riemann Zero Operators]{%
  Spectral Exclusion Results for Operators Whose Eigenvalue
  Multiset Would Equal the Riemann Zero Ordinates}

\author{Lin Tao}
\date{2026}

\begin{document}

\begin{abstract}
We prove two spectral exclusion results showing that no operator in
a natural geometric class can have its eigenvalue multiset equal to
$\{\gamma_n\}$, the positive ordinates of the nontrivial zeros of
the Riemann zeta function.

\textbf{Theorem~D.}  No positive classical elliptic
pseudodifferential operator of order $m>0$ on any compact smooth
manifold has eigenvalue multiset $\{\gamma_n\}$: the Weyl counting
law forces a pure-power leading term in the heat trace, while the
zero ordinate heat sum $Z_\zeta(t)\sim\frac{1}{2\pi}t^{-1}\log(1/t)$
has an unavoidable logarithmic factor.

\textbf{Theorem~D$'$.}  The exclusion extends to the class~$\Csub$
of operators with log-polyhomogeneous sub-principal symbols,
via a min-max squeeze argument.

Both results are unconditional; neither assumes nor concludes the
Riemann Hypothesis.
\end{abstract}

\maketitle
\tableofcontents

%-----------------------------------------------------------------------
\section{Introduction}
\label{sec:intro}
%-----------------------------------------------------------------------

\subsection*{Setup and motivation}

The nontrivial zeros of $\zeta(s)$ in the upper half-plane may be
written as $\rho_n = \beta_n + i\gamma_n$ with $\beta_n\in(0,1)$ and
$\gamma_n>0$ (listed with multiplicity, $0<\gamma_1\le\gamma_2\le\cdots$).
The ordinates $\gamma_n$ are unconditionally real; the Riemann
Hypothesis asserts that $\beta_n=\tfrac{1}{2}$ for all $n$.

A programme going back to Hilbert and P\'olya asks whether a
self-adjoint operator $H$ on a Hilbert space with a natural geometric
or analytic origin can be exhibited with eigenvalue multiset
$\{\gamma_n\}$ (eigenvalues listed with multiplicity).
Even granted a formal construction, this spectral matching question
is a prerequisite for any programme linking $H$ to RH.
The present paper asks which naturally defined operator classes are
\emph{provably incapable} of achieving this match.

\subsection*{Results}

\begin{enumerate}[label=(\textbf{\Alph*})]
  \setcounter{enumi}{3}
  \item[\textbf{D}] (Section~\ref{sec:elliptic}).  Classical elliptic
    pseudodifferential operators on compact manifolds are excluded
    by a heat-trace comparison.
  \item[\textbf{D$'$}] (Section~\ref{sec:logpoly}).  The exclusion
    extends to the class $\Csub$ of operators with
    log-polyhomogeneous sub-principal symbols.
\end{enumerate}

Both results are RH-free; they establish structural constraints
independent of the truth or falsity of RH.

\subsection*{Related work}

Watson--Valentinuzzi~\cite{Watson2026} prove a Weyl-law obstruction
for compact elliptic \emph{differential} operators; Theorem~D gives
the same obstruction for classical elliptic pseudodifferential
operators of any order $m>0$ via the pseudodifferential Weyl law
(H\"ormander~\cite{Hormander1968}).
Theorem~D$'$ extends the result to the wider log-polyhomogeneous
symbol class of Lesch~\cite{Lesch1999}.
Connes~\cite{Connes1999} places zeta zeros in an adelic spectral
setting.  Connes--Consani--Moscovici~\cite{CCM2025} construct
scaled Dirac-type operators $D_{\!\log}^{\lambda,N}$ on an adelic
space and establish a regularized determinant identity involving a
scaled entire function; showing that the resulting determinants
converge to the Riemann $\Xi$ is the open step in their programme.
Theorems~D and~D$'$ do not obstruct the CCM approach, since their
operators belong to a different class than the CCM construction.

%-----------------------------------------------------------------------
\section{Setup and Notation}
\label{sec:setup}
%-----------------------------------------------------------------------

\subsection{Heat traces and counting functions}

Let $H$ be a non-negative self-adjoint operator with compact resolvent
on a separable Hilbert space, with eigenvalues
$0\le e_1\le e_2\le\cdots$ listed with multiplicity.
We refer to this listed sequence as the \emph{eigenvalue multiset}
of $H$ and write it $\operatorname{spec}_{\times}(H)$.
The \emph{heat trace} and \emph{spectral counting function} are
\[
  Z_H(t)=\operatorname{Tr}(e^{-tH})=\sum_{n=1}^\infty e^{-e_n t},
  \quad t>0; \qquad
  N_H(\Lambda)=\#\{n:e_n\le\Lambda\}.
\]

\begin{theorem}[Karamata~\cite{BGT1987}, Theorem~1.7.1]
\label{thm:karamata}
For $\alpha>0$ and $C>0$,
\[
  N_H(\Lambda)\sim C\Lambda^\alpha
  \;\iff\;
  Z_H(t)\sim C\,\Gamma(\alpha+1)\,t^{-\alpha}
  \quad(t\to 0^+).
\]
\end{theorem}

\subsection{The zero ordinate heat sum}

Write the nontrivial zeros in the upper half-plane as
$\rho_n=\beta_n+i\gamma_n$ ($\beta_n\in(0,1)$, $\gamma_n>0$, listed
with multiplicity).  Define $Z_\zeta(t)=\sum_{n=1}^\infty e^{-\gamma_n t}$.
The Riemann--von Mangoldt formula
(Titchmarsh--Heath-Brown~\cite{Titchmarsh1986}, Theorem~9.4) gives
\[
  N_\zeta(T)\;:=\;\#\{n:\gamma_n\le T\}
  \;\sim\;\frac{T\log T}{2\pi}
  \quad(T\to+\infty),
\]
unconditionally.  Applying the Tauberian form of Karamata's theorem
for logarithmically varying functions (\cite{BGT1987}, \S1.7) yields
\begin{equation}
  \label{eq:zetaheat}
  Z_\zeta(t)\sim\frac{1}{2\pi}\,t^{-1}\log(1/t)
  \quad(t\to 0^+).
\end{equation}

\subsection{Manifold conventions}

All manifolds are closed (compact, without boundary) and smooth unless
otherwise stated.  A \emph{classical elliptic pseudodifferential
operator} of order $m>0$ on a $d$-dimensional manifold $M$ is a
positive self-adjoint operator $H$ with compact resolvent whose
complete symbol is polyhomogeneous of order $m$ and whose principal
symbol $h_m(x,\xi)$ satisfies $h_m(x,\xi)\ge c|\xi|^m$ uniformly
for $|\xi|\ge 1$.

%-----------------------------------------------------------------------
\section{Classical Elliptic Operators Are Excluded}
\label{sec:elliptic}
%-----------------------------------------------------------------------

\begin{theorem}[D]
\label{thm:D}
Let $M$ be a closed smooth $d$-manifold and $H$ a positive classical
elliptic pseudodifferential operator of order $m>0$ on $M$ with
compact resolvent.  Then $\operatorname{spec}_{\times}(H)\ne\{\gamma_n\}$.
\end{theorem}

\begin{proof}
By the pseudodifferential Weyl law
(H\"ormander~\cite{Hormander1968}; for order $m>0$ reduce to order $1$
via $H^{1/m}$, which is a classical elliptic $\Psi$DO of order 1 by the
symbolic calculus, with $N_H(\Lambda)=N_{H^{1/m}}(\Lambda^{1/m})$),
\[
  N_H(\Lambda)\sim C_H\,\Lambda^{d/m}
  \quad(\Lambda\to+\infty),
\]
where
\[
  C_H=(2\pi)^{-d}\int_{T^*M}\mathbf{1}_{h_m(x,\xi)\le 1}\,d\xi\,dx\;>\;0.
\]
By Theorem~\ref{thm:karamata},
\begin{equation}
  \label{eq:classheat}
  Z_H(t)\sim C_H\,\Gamma(d/m+1)\,t^{-d/m}
  \quad(t\to 0^+).
\end{equation}
In particular, the leading term of $Z_H(t)$ as $t\to 0^+$ is a pure
power $t^{-d/m}$.

Suppose for contradiction that $\operatorname{spec}_{\times}(H)=\{\gamma_n\}$.
Then $Z_H(t)=Z_\zeta(t)$ for all $t>0$, so
by~\eqref{eq:zetaheat} and~\eqref{eq:classheat}:
\[
  C_H\,\Gamma(d/m+1)\,t^{-d/m}
  \;\sim\;
  \frac{1}{2\pi}\,t^{-1}\log(1/t).
\]
Let $\alpha=d/m>0$.  As $t\to 0^+$, the ratio of the two sides is
\[
  \frac{C_H\Gamma(\alpha+1)\,t^{-\alpha}}{(1/2\pi)t^{-1}\log(1/t)}
  \;=\;
  2\pi C_H\Gamma(\alpha+1)\,
  \frac{t^{1-\alpha}}{\log(1/t)}.
\]
\begin{itemize}
\item If $\alpha<1$: $t^{1-\alpha}\to 0$ while $\log(1/t)\to+\infty$,
  so the ratio $\to 0$.
\item If $\alpha>1$: $t^{1-\alpha}\to+\infty$ faster than
  $\log(1/t)\to+\infty$, so the ratio $\to+\infty$.
\item If $\alpha=1$: the ratio becomes
  $2\pi C_H\Gamma(2)/\log(1/t)\to 0$ as $t\to 0^+$.
\end{itemize}
In all three cases the two sides have different asymptotic behavior,
contradicting $Z_H(t)=Z_\zeta(t)$.
\end{proof}

\begin{remark}
\label{rem:bgvgilkey}
The heat-kernel results of Berline--Getzler--Vergne~\cite{BGV2004}
(Theorem~2.30) and Gilkey~\cite{Gilkey1995} (\S1.7) cover
Laplace-type and differential operators; those results suffice for
a restricted version of Theorem~D but do not cover general
classical elliptic pseudo\-differential operators.
\end{remark}

\begin{remark}[Watson--Valentinuzzi]
Watson and Valentinuzzi~\cite{Watson2026} prove a Weyl-law obstruction
for compact elliptic differential operators.  Theorem~D extends
their result to classical elliptic pseudo\-differential operators
of any order $m>0$.
\end{remark}

\subsection*{Escape routes from Theorem~D}

The proof uses the classical Weyl law, which requires: (i) $M$ compact,
(ii) the operator elliptic, (iii) the symbol polyhomogeneous.
Genuine escape routes are:
\begin{enumerate}[label=(\arabic*)]
\item \emph{Non-elliptic operators}---principal symbol vanishing on a
  characteristic set can produce non-pure-power counting asymptotics
  (see the counterexample in Section~\ref{sec:logpoly}).
\item \emph{Operators on non-compact manifolds.}
\item \emph{Sub-classical symbol growth}, e.g.\ $|\xi|/\log|\xi|$-type
  symbols; see the discussion in Section~\ref{sec:logpoly}.
\end{enumerate}

%-----------------------------------------------------------------------
\section{Extension to Log-Polyhomogeneous Symbols}
\label{sec:logpoly}
%-----------------------------------------------------------------------

\subsection{The class \texorpdfstring{$\Csub$}{C-logpoly}}

\begin{definition}
\label{def:csub}
The class $\Csub$ consists of scalar self-adjoint operators $H\ge -C_0$
(some $C_0\ge 0$) with compact resolvent on a closed $d$-manifold $M$
whose complete symbol has the asymptotic expansion
\[
  \sigma(H)(x,\xi)
  \sim
  h_m(x,\xi)
  +
  \sum_{j\ge 1}\sum_{\ell=0}^{K}
    h_{m-j,\ell}(x,\xi)(\log|\xi|)^\ell
  \quad(|\xi|\to\infty),
\]
where each $h_{m-j,\ell}$ is homogeneous of degree $m-j$, and the
principal symbol $h_m$ satisfies $h_m(x,\xi)\ge c|\xi|^m$ for
$|\xi|\ge 1$ (positive classical elliptic principal part).
\end{definition}

The log-polyhomogeneous calculus is developed in
Grubb--Seeley~\cite{GrubbSeeley1995} and Lesch~\cite{Lesch1999}.

\subsection{Form boundedness of the log correction}

The sub-principal log terms constitute an operator $Q=H-P$, where $P$
is the classical elliptic operator with principal symbol $h_m$ and no
log corrections.

For any $\varepsilon>0$, each term
$h_{m-j,\ell}(x,\xi)(\log|\xi|)^\ell$ with $j\ge 1$ satisfies
$|h_{m-j,\ell}(x,\xi)(\log|\xi|)^\ell|\le C|\xi|^{m-1+\varepsilon}$
for $|\xi|\ge 1$ (since $|\log|\xi||=O(|\xi|^\varepsilon)$
for every $\varepsilon>0$).  Thus $Q\in\Psi^{m-1+\varepsilon}(M)$ for
every $\varepsilon>0$.

Since $P$ is a positive elliptic operator of order $m$ on the compact
manifold $M$, the Sobolev scale satisfies
$\|u\|_{H^{m/2}(M)}^2\asymp\langle Pu,u\rangle+\|u\|_{L^2}^2$.
For $u\in H^{m/2}(M)$ and any $\varepsilon\in(0,1)$:
\[
  |\langle Qu,u\rangle|
  \le \|Qu\|_{H^{-m/2}}\,\|u\|_{H^{m/2}}.
\]
Since $Q\in\Psi^{m-1+\varepsilon}$, the map
$Q:H^{m/2}\to H^{1/2-\varepsilon}$ is bounded, and
$H^{1/2-\varepsilon}\hookrightarrow H^{-m/2}$ (as $1/2-\varepsilon>-m/2$
for $m<1+2\varepsilon$, and for larger $m$ one picks $\varepsilon$ small).
More directly, the Sobolev interpolation inequality on $M$ gives:
for every $\delta>0$ there exists $C_\delta>0$ such that
\[
  \|u\|_{H^{(m-1+\varepsilon)/2}(M)}^2
  \;\le\;
  \delta\,\|u\|_{H^{m/2}(M)}^2+C_\delta\,\|u\|_{L^2(M)}^2,
\]
since $(m-1+\varepsilon)/2<m/2$ for $\varepsilon<1$.  Using
$|\langle Qu,u\rangle|\le C'\|u\|_{H^{(m-1+\varepsilon)/2}}^2$
(Calderón--Vaillancourt: $Q$ maps $H^{(m-1+\varepsilon)/2}$ to
its dual $H^{-(m-1+\varepsilon)/2}$, and the form is $L^2$ pairing),
we obtain: for every $\delta>0$,
\begin{equation}
  \label{eq:formbdd}
  |\langle Qu,u\rangle|
  \;\le\;
  \delta\langle Pu,u\rangle+C_\delta\|u\|^2
  \qquad\forall\,u\in H^{m/2}(M).
\end{equation}
Equivalently, as quadratic forms on $H^{m/2}(M)$:
\begin{equation}
  \label{eq:formineq}
  (1-\delta)P-C_\delta\;\le\; H\;\le\;(1+\delta)P+C_\delta.
\end{equation}

\begin{theorem}[D$'$]
\label{thm:Dprime}
Any $H\in\Csub$ satisfies
$N_H(\Lambda)\sim C_H\Lambda^{d/m}$ as $\Lambda\to\infty$,
where $C_H>0$ is the Weyl constant of the principal symbol $h_m$.
In particular, $\operatorname{spec}_{\times}(H)\ne\{\gamma_n\}$.
\end{theorem}

\begin{proof}
Write $a=C_0+1$.  Replace $H$ by $\tilde H=H+a$, which satisfies
$\tilde H\ge 1>0$.  The counting functions are related by
$N_{\tilde H}(\Lambda)=N_H(\Lambda-a)$, so
$N_H(\Lambda)\sim C_H\Lambda^{d/m}$ is equivalent to
$N_{\tilde H}(\Lambda)\sim C_H\Lambda^{d/m}$; the shift $a$ does not
affect the leading Weyl asymptotics.  Denote $\tilde H$ by $H$ for
the rest of the proof.

The comparison operator $P$ satisfies $P\ge c>0$ and the
form inequality~\eqref{eq:formineq} holds for $H$.

By the min-max (Courant--Fischer) principle, if $A\le B$ as quadratic
forms on a common domain, then $\lambda_k(A)\le\lambda_k(B)$ for
every $k$.  Applying this to~\eqref{eq:formineq}:
\begin{equation}
  \label{eq:minmax}
  N_P\!\Bigl(\frac{\Lambda-C_\delta}{1+\delta}\Bigr)
  \;\le\;
  N_H(\Lambda)
  \;\le\;
  N_P\!\Bigl(\frac{\Lambda+C_\delta}{1-\delta}\Bigr),
  \qquad 0<\delta<1.
\end{equation}
Justification: from $H\le(1+\delta)P+C_\delta$, we have
$\lambda_k(H)\le(1+\delta)\lambda_k(P)+C_\delta$, so
$\lambda_k(P)\le(\Lambda-C_\delta)/(1+\delta)$
\emph{implies} $\lambda_k(H)\le\Lambda$.
Hence $\{k:\lambda_k(P)\le(\Lambda-C_\delta)/(1+\delta)\}\subseteq
\{k:\lambda_k(H)\le\Lambda\}$, giving the lower bound.
From $(1-\delta)P-C_\delta\le H$, we have
$\lambda_k(H)\ge(1-\delta)\lambda_k(P)-C_\delta$, so
$\lambda_k(H)\le\Lambda$ \emph{implies}
$\lambda_k(P)\le(\Lambda+C_\delta)/(1-\delta)$.
Hence $\{k:\lambda_k(H)\le\Lambda\}\subseteq
\{k:\lambda_k(P)\le(\Lambda+C_\delta)/(1-\delta)\}$,
giving the upper bound.

Since $N_P(\Lambda)\sim C_H\Lambda^{d/m}$ (Theorem~D, applied to $P$),
dividing~\eqref{eq:minmax} by $\Lambda^{d/m}$ and letting
$\Lambda\to+\infty$ gives
\[
  C_H(1+\delta)^{-d/m}
  \;\le\;\liminf_{\Lambda\to\infty}\frac{N_H(\Lambda)}{\Lambda^{d/m}}
  \;\le\;\limsup_{\Lambda\to\infty}\frac{N_H(\Lambda)}{\Lambda^{d/m}}
  \;\le\;C_H(1-\delta)^{-d/m}.
\]
Since $\delta>0$ is arbitrary, $N_H(\Lambda)\sim C_H\Lambda^{d/m}$.
By Theorem~\ref{thm:karamata}, $Z_H(t)\sim C_H\Gamma(d/m+1)t^{-d/m}$,
and the exclusion follows by the same argument as Theorem~D.
\end{proof}

\begin{remark}
Lesch~\cite{Lesch1999} Theorem~3.7 establishes the resolvent kernel
expansion for log-polyhomogeneous operators, from which eq.~(3.18)
of that paper derives heat-trace asymptotics for
$\operatorname{Tr}(Ae^{-tP})$ with log-polyhomogeneous $A$ and
classical elliptic $P$.  This does not directly give $Z_H(t)$
asymptotics when $H$ itself is log-polyhomogeneous.
Our proof avoids this gap by using the form-boundedness
comparison~\eqref{eq:formineq} and the classical Weyl law for $P$,
together with the min-max principle.
\end{remark}

\subsection{A genuine non-elliptic escape}

The following example shows that log-enhanced heat traces \emph{can}
occur once ellipticity is dropped.

\begin{remark}[Non-elliptic counterexample]
\label{rem:counterex}
Let $\mathbb{T}^4=\mathbb{T}^2_x\times\mathbb{T}^2_y$ and set
$H_{\mathrm{prod}}=(I-\Delta_x)(I-\Delta_y)$.
The eigenvalues are $(1+|p|^2)(1+|q|^2)$ for $(p,q)\in\ZZ^2\times\ZZ^2$
(using integer Fourier modes for $\mathbb{T}^2=\RR^2/(2\pi\ZZ)^2$).
We claim $N_{H_{\mathrm{prod}}}(\Lambda)\sim\pi^2\Lambda\log\Lambda$.

For each $p\in\ZZ^2$, the number of $q\in\ZZ^2$ with
$(1+|q|^2)\le\Lambda/(1+|p|^2)$ is asymptotically
$\pi\Lambda/(1+|p|^2)$ (area of the disk of radius
$\sqrt{\Lambda/(1+|p|^2)}$).
Summing over $p$ with $1+|p|^2\le\Lambda$:
\begin{align*}
  N_{H_{\mathrm{prod}}}(\Lambda)
  &\sim
  \pi\Lambda\sum_{\substack{p\in\ZZ^2\\|p|^2\le\Lambda-1}}
  \frac{1}{1+|p|^2}.
\end{align*}
The sum $\sum_{p\in\ZZ^2,\,|p|^2\le R}(1+|p|^2)^{-1}\sim\pi\log R$
as $R\to\infty$ (the two-dimensional lattice harmonic series,
computed by grouping terms at $|p|^2=k$ and using
$\#\{p\in\ZZ^2:|p|^2=k\}=O(k^\varepsilon)$ with a sum converging
logarithmically).
Taking $R=\Lambda-1\sim\Lambda$ gives
$N_{H_{\mathrm{prod}}}(\Lambda)\sim\pi\Lambda\cdot\pi\log\Lambda
=\pi^2\Lambda\log\Lambda$.

By Karamata,
$Z_{H_{\mathrm{prod}}}(t)\sim\pi^2 t^{-1}\log(1/t)$ as $t\to 0^+$.
This operator is \emph{not} elliptic on $\mathbb{T}^4$: its principal
symbol $|\xi|^2|\eta|^2$ vanishes on the characteristic set
$\{\xi=0\}\cup\{\eta=0\}$, which is non-empty.  Hence logarithmically
enhanced heat traces are achievable outside the elliptic class.
\end{remark}

\subsection*{The genuine escape class}

To match $Z_\zeta(t)$ one needs $N_H(T)\sim CT\log T$.
A heuristic that achieves this counting law is the Lambert-W corrected
symbol on the complete circle $S^1$ (both positive and negative
Fourier modes):
\[
  g(r)=\frac{4\pi r}{W(2r/e)},
\]
where $W$ is the Lambert-W function ($W(x)e^{W(x)}=x$).  The factor
$4\pi$ accounts for the double-sided spectrum $\{\pm n\}_{n\ge 1}$
on $S^1$: if $g^{-1}(T)\sim\frac{T}{4\pi}\log(T/2)$,
then $N(T)\approx 2g^{-1}(T)\sim\frac{T}{2\pi}\log T$, matching
the first term of $N_\zeta(T)$.  Whether a self-adjoint realization
with eigenvalue multiset exactly $\{\gamma_n\}$ can be constructed
from such symbols remains open.

%-----------------------------------------------------------------------
\section{Discussion}
\label{sec:discussion}
%-----------------------------------------------------------------------

Theorems~D and~D$'$ close off the classical elliptic and
log-poly\-homo\-geneous pseudo\-differential regime on compact manifolds.
The obstruction is the Weyl law: classical elliptic operators have
a pure-power leading term in the heat trace, incompatible with the
log-enhanced $Z_\zeta(t)\sim(1/2\pi)t^{-1}\log(1/t)$.
Theorem~D$'$ shows this obstruction is stable under
log-poly\-homo\-geneous lower-order perturbations.
The $\mathbb{T}^4$ counterexample (Remark~\ref{rem:counterex})
confirms that log-enhancement is achievable outside the elliptic
class, so the ellipticity hypothesis in Theorem~D is sharp.

Open questions:
\begin{enumerate}[label=(\arabic*)]
\item Can the CCM normalization step be completed, giving a
  spectral realization of $\Xi$ from the $D_{\!\log}^{\lambda,N}$
  family?
\item Can a self-adjoint operator with $|\xi|/\log|\xi|$-type symbol
  on some manifold realize $\{\gamma_n\}$ as its eigenvalue multiset?
\item What is the minimal analytic condition on a symbol class that
  permits $N_H(T)\sim CT\log T$ while retaining a functional calculus?
\end{enumerate}

\begin{thebibliography}{99}

\bibitem{BGT1987}
N.~H. Bingham, C.~M. Goldie, and J.~L. Teugels,
\emph{Regular Variation},
Cambridge University Press, Cambridge, 1987.

\bibitem{BGV2004}
N.~Berline, E.~Getzler, and M.~Vergne,
\emph{Heat Kernels and Dirac Operators},
Grundlehren der Mathematischen Wissenschaften, vol.~298,
Springer-Verlag, Berlin, 2004.

\bibitem{Connes1999}
A.~Connes,
Trace formula in noncommutative geometry and the zeros of the Riemann
zeta function,
\emph{Selecta Math.}\ \textbf{5} (1999), 29--106.

\bibitem{CCM2025}
A.~Connes, C.~Consani, and H.~Moscovici,
Zeta spectral triples,
\emph{arXiv:2511.22755} (2025).

\bibitem{GrubbSeeley1995}
G.~Grubb and R.~T. Seeley,
Weakly parametric pseudodifferential operators and
Atiyah--Patodi--Singer boundary problems,
\emph{Invent.\ Math.}\ \textbf{121} (1995), no.~3, 481--529.

\bibitem{Gilkey1995}
P.~B. Gilkey,
\emph{Invariance Theory, the Heat Equation, and the
Atiyah--Singer Index Theorem},
2nd ed., CRC Press, Boca Raton, FL, 1995.

\bibitem{Hormander1968}
L.~H\"ormander,
The spectral function of an elliptic operator,
\emph{Acta Math.}\ \textbf{121} (1968), 193--218.

\bibitem{Lesch1999}
M.~Lesch,
On the noncommutative residue for pseudodifferential operators with
log-polyhomogeneous symbols,
\emph{Ann.\ Global Anal.\ Geom.}\ \textbf{17} (1999), no.~2, 151--187.

\bibitem{Titchmarsh1986}
E.~C. Titchmarsh,
\emph{The Theory of the Riemann Zeta-Function},
2nd ed., revised by D.~R. Heath-Brown,
Oxford University Press, Oxford, 1986.

\bibitem{Watson2026}
D.~F. Watson and T.~Valentinuzzi,
Spectral-dimension obstructions for operators with superlinear
counting laws,
\emph{Bull.\ Sci.\ Math.}\ (2026), arXiv:2604.00052.

\end{thebibliography}

\end{document}
